\documentclass{article}
\usepackage[T1]{fontenc}
\usepackage{lmodern}
\usepackage{graphicx}
\usepackage{amsmath,amssymb}
\usepackage[a4paper,margin=2cm]{geometry}
\usepackage[absolute,overlay]{textpos}
\setlength{\TPHorizModule}{1mm}
\setlength{\TPVertModule}{1mm}
\usepackage[indonesian]{babel}
\usepackage{hyperref}
\setlength{\emergencystretch}{2em}
% unit: Halaman 53

\begin{document}

% === Halaman 53 (llm) ===

\section*{1. Confusion}

\begin{itemize}
    \item Prinsip ini menyembunyikan hubungan statistik yang ada diantara plainteks, cipherteks, dan kunci.
    \item Prinsip \textit{confusion} membuat kriptanalis frustasi untuk mencari pola-pola statistik yang muncul di dalam cipherteks.
    \item \textit{Confusion} dapat direalisasikan dengan menggunakan teknik substitusi yang nirlanjar (\textit{non-linear}), biasanya dengan operasi \textit{lookup table}.
\end{itemize}

\end{document}
